% mode: ini
% Aimed at append_discretionary and build_discretionary: \- with hyphen_char in and out of range and a missing one; each of the
% three parts, each allowed item, each improper one with the display of what was deleted; "Illegal math \discretionary"; and
% the replacement's length: replace_count, 255 and 256 items against "Discretionary list is too long", where the limit is
% WEB's max_quarterword.  Also scanning: a part that is not in braces, and a \par in a part.
\input prelude
\let\bgroup={ \let\egroup=}
\output={\setbox0\box255 \deadcycles=0 }
\def\reset{\noindent\par}
\errorcontextlines=0
\tracinglostchars=1
\def\b#1{\setbox0\hbox{#1}\showbox0 }
% \- and hyphen_char.
\b{\-}\b{a\-b}\noindent\-\showlists\par\reset
$\-\showlists$\reset
\hyphenchar\tenrm=-1 \b{\-}
\hyphenchar\tenrm=255 \b{\-}
\hyphenchar\tenrm=256 \b{\-}
\hyphenchar\tenrm=0 \b{\-}
\hyphenchar\tenrm=`- \nullfont\b{\-}\tenrm
\-\showlists\par\reset
\discretionary{}{}{}\showlists\par\reset
% Each part with each item.
\b{\discretionary{a}{b}{c}}
\b{\discretionary{fi}{\kern1pt}{\hbox{}}}
\b{\discretionary{\vrule}{\vbox{}}{\vtop{}}}
\b{\discretionary{\accent"7F a}{}{}}
\b{\discretionary{}{}{ab\kern1pt\vrule\hbox{}fi}}
\b{x\discretionary{}{}{ab\kern1pt}y}
\b{\discretionary{\/}{}{}}
\b{\discretionary\bgroup a\egroup\bgroup b\egroup\bgroup c\egroup}
\b{\discretionary{}{}{{a}{b}}}
\b{\discretionary{\begingroup a\endgroup}{}{}}
\b{\discretionary{\setbox1\hbox{}\box1}{\count0=1 }{}}
% Improper items, each kind, in each part, and with something before and after to be flushed.
\def\try#1{\b{\discretionary{#1}{}{}}\b{\discretionary{}{#1}{}}\b{\discretionary{}{}{#1}}\b{\discretionary{}{}{x#1y}}}
\try{\penalty1 }
\try{\hskip1pt}
\try{\mark{m}}
\try{\insert100{}}
\try{\vadjust{}}
\try{\discretionary{}{}{}}
\try{\-}
\try{$$}
\try{\special{s}}
\try{\write1{w}}
\try{ }
% Math formula and glue left by \raise.  Now in math mode.
$\discretionary{a}{b}{}\showlists$\reset
$\discretionary{a}{b}{c}\showlists$\reset
$\discretionary{}{}{\kern1pt}\showlists$\reset
$\discretionary{}{}{\hbox{}}\showlists$\reset
$\discretionary{}{}{\penalty0 c}\showlists$\reset
$$\discretionary{a}{b}{c}\showlists$$\reset
$$\discretionary{}{}{}\showlists$$\reset
% In paragraphs, for the \tail bookkeeping: replacement goes after the node, and the tail moves to its end.
\noindent x\discretionary{}{}{ab}y\showlists\par\reset
\noindent x\discretionary{}{}{}\showlists\par\reset
\noindent x\discretionary{}{}{ab}\unkern\showlists\par\reset
% A paragraph in a part: the group is restricted horizontal.
\noindent\discretionary{\par}{\par}{\par}\showlists\par\reset
\noindent\discretionary{\vskip1pt}{}{}\showlists\par\reset
% Scanning.
\noindent\discretionary x{}{}\showlists\par\reset
\noindent\discretionary{}x{}\showlists\par\reset
\noindent\discretionary{}{}x\showlists\par\reset
\noindent\discretionary\relax\relax{}\relax{}\relax{}\showlists\par\reset
% Length.  Set up 254, 255, 256, and 257 items.
\def\ten{aaaaaaaaaa}\def\hundred{\ten\ten\ten\ten\ten\ten\ten\ten\ten\ten}
\def\ccl{\hundred\hundred\ten\ten\ten\ten\ten}
\hyphenchar\tenrm=-1
\setbox0\hbox{\discretionary{}{}{\ccl aaaa}\showlists}\showboxbreadth=3 \showbox0 \showboxbreadth=100
\setbox0\hbox{\discretionary{}{}{\ccl aaaaa}\showlists}\showboxbreadth=3 \showbox0 \showboxbreadth=100
\showboxbreadth=3 \showboxdepth=1
\setbox0\hbox{\discretionary{}{}{\ccl aaaaaa}x\showlists}\showbox0
\setbox0\hbox{\discretionary{}{}{\ccl aaaaaaa}x\showlists}\showbox0
\showboxbreadth=100 \showboxdepth=100
% A long list in the math part: "Illegal" comes first.
$\discretionary{}{}{\ccl aaaaaa}\showlists$\reset
\hyphenchar\tenrm=`-
\showlists
\end
